\documentclass[10pt]{article}
\usepackage[margin=0.85in]{geometry}
\usepackage[T1]{fontenc}
\usepackage{lmodern}
\usepackage{amsmath,amssymb,amsthm}
\usepackage{booktabs}
\usepackage{microtype}
\usepackage[hidelinks]{hyperref}
\newtheorem{theorem}{Theorem}
\newtheorem{proposition}{Proposition}
\newcommand{\Z}{\mathbb Z}
\newcommand{\Hom}{\operatorname{Hom}}
\newcommand{\im}{\operatorname{im}}
\title{Intersection Geometry of Complements in a Restricted Mark Cokernel}
\author{Anonymous}
\date{}
\begin{document}
\maketitle

\begin{abstract}
An exact restricted mark cokernel was previously split as $C=D\oplus K$, with
$D\cong\mathbb Z/8\oplus(\mathbb Z/2)^2$ and exactly $2^{41}$ complements.
We determine the intersection geometry of the entire complement family.
Relative to any fixed complement, the possible intersection indices are
$1,2,4,8,16,32$, and we give the exact multiplicity of every index and every
intersection-quotient type.  The common intersection of all complements is
$8C\cong\mathbb Z/3\oplus\mathbb Z/18$, whereas their generated span is all
of $C$.  Finally, we compute the orders of the sixteen named classes in this
common core and find exactly twenty-five generating triples, all containing
the ninth class.  Every claim remains on the frozen sixteen-type support.
\end{abstract}

\section{Graph model and scope}
Fix the C69 decomposition $C=D\oplus K$.  In primary exponent notation,
\[
 D:(3,1,1),\qquad K_{(2)}:(4,2,2,2,1^8),\qquad K_{(3)}:(2,1).
\]
Every complement of $D$ is uniquely the graph
\[
 \Gamma_f=\{(f(k),k):k\in K\},\qquad f\in\Hom(K,D).
\]
The parameter group is
\[
 \Hom(K,D)\cong(\Z/2)^{32}\oplus(\Z/4)^3\oplus\Z/8,
 \qquad |\Hom(K,D)|=2^{41}.
\]
For two parameters $f,g$ one has
\begin{equation}\label{eq:graph}
 \Gamma_f\cap\Gamma_g\cong\ker(f-g),\qquad
 K/\ker(f-g)\cong\im(f-g).
\end{equation}
Thus all intersection statistics are translation invariant in the affine
parameter group.  Throughout, ``intersection index'' means
\[
 [\Gamma_f:\Gamma_f\cap\Gamma_g]=[K:\ker(f-g)]=|\im(f-g)|,
\]
not the index of the intersection in $C$.

The scope firewall is
\[
\texttt{NO\_BAD\_EULER\_OR\_ROOT\_NUMBER}.
\]
We do not claim a canonical complement, a complete classification of the
abstract kernel types, a full table of marks or Burnside ring, arithmetic or
local data, Euler factors, root numbers, automorphy, or a Hilbert--Polya
operator.

\section{Exact quotient-type distribution}
For a subgroup type $H\leq D$, the number of maps $K\to D$ with image of type
$H$ is
\begin{equation}\label{eq:imagecount}
 N_D(H)\,E_K(H),
\end{equation}
where $N_D(H)$ counts subgroups of $D$ of type $H$ and $E_K(H)$ counts
epimorphisms from $K$ onto one fixed copy of $H$.  The producer obtains
$E_K(H)$ by strict inversion on the concrete 38-element subgroup poset of
$D$.  Independently, the checker uses
\[
 E_K(H)=N_K(H)|\operatorname{Aut}(H)|,
\]
because finite duality sends epimorphisms $K\to H$ to injections
$H^\vee\to K^\vee$, and these groups are abstractly self-dual.  Birkhoff's
subgroup formula evaluates both $N_D(H)$ and $N_K(H)$.

\begin{table}[htbp]
\centering
\small
\begin{tabular}{lrrr}
\toprule
$H\cong K/\ker f$ & $N_D(H)$ & $E_K(H)$ & maps with image type $H$\\
\midrule
$0$ & 1 & 1 & 1\\
$\Z/2$ & 7 & 4095 & 28665\\
$\Z/4$ & 4 & 61440 & 245760\\
$(\Z/2)^2$ & 7 & 16764930 & 117354510\\
$\Z/8$ & 4 & 65536 & 262144\\
$\Z/2\oplus\Z/4$ & 6 & 251535360 & 1509212160\\
$(\Z/2)^3$ & 1 & 68602093560 & 68602093560\\
$\Z/2\oplus\Z/8$ & 6 & 268304384 & 1609826304\\
$(\Z/2)^2\oplus\Z/4$ & 1 & 1029282693120 & 1029282693120\\
$\Z/8\oplus(\Z/2)^2$ & 1 & 1097901539328 & 1097901539328\\
\bottomrule
\end{tabular}
\caption{Exact image, hence intersection-quotient, distribution from every
fixed complement.}
\label{tab:types}
\end{table}

\begin{theorem}[Intersection spectrum]
From every fixed complement, the numbers of complements at intersection
indices $1,2,4,8,16,32$ are respectively
\[
 1,\ 28665,\ 117600270,\ 70111567864,\
 1030892519424,\ 1097901539328.
\]
Their sum is $2^{41}$.  Table~\ref{tab:types} refines these six values into
all ten possible quotient types.
\end{theorem}

\begin{proof}
Equation~\eqref{eq:graph} makes the index equal to $|\im(f-g)|$.  Summing
Equation~\eqref{eq:imagecount} over subgroup types of each order gives the
displayed values.  Translation by $f$ is a bijection of $\Hom(K,D)$, so the
same row occurs at every fixed complement.  Multiplying by $2^{41}$ gives
ordered-pair counts; halving the off-diagonal rows gives unordered-pair
counts, all recorded in the exact evidence.
\end{proof}

\section{The common core and total span}
An element common to every graph must first lie in $\Gamma_0$, hence has the
form $(0,k)$, and it belongs to all graphs precisely when every homomorphism
$K\to D$ kills $k$.  Since $D$ has exponent eight, $8K$ is contained in this
universal kernel.  Conversely, the cyclic $\Z/8$ summand of $D$ separates
every class outside $8K$ on the dyadic primary part.  The odd-primary part is
killed by every map to $D$ and multiplication by eight is invertible there.
Therefore the universal kernel is exactly $8K$.

\begin{theorem}[Core and span]
The complement family satisfies
\[
 \bigcap_f\Gamma_f=8C\cong\Z/3\oplus\Z/18,
 \qquad |8C|=54,
\]
and
\[
 \left\langle\Gamma_f:f\in\Hom(K,D)\right\rangle=C.
\]
Indeed, from every fixed complement there are $1097901539328$ other
parameters whose difference is surjective, and each such pair already
generates $C$.
\end{theorem}

\begin{proof}
Because $8D=0$, the preceding universal-kernel argument identifies the
intersection with $8K=8C$.  Multiplication by eight on the primary factors of
$K$ leaves one $\Z/2$, together with $\Z/3\oplus\Z/9$; combining coprime
factors gives $\Z/3\oplus\Z/18$.  If $f-g$ is surjective, differences of
elements of $\Gamma_f$ and $\Gamma_g$ generate $D$, while either graph maps
onto $K$.  The last row of Table~\ref{tab:types} counts all such differences.
\end{proof}

\section{Named generators of the core}
Let $e_1,\ldots,e_{16}$ be the named coordinate classes in the frozen mark
presentation.  Exact rational residues, or independently Hermite lattices,
give
\[
 \bigl(|8[e_j]|\bigr)_{j=1}^{16}
 =(9,3,3,3,1,1,9,3,2,1,3,3,1,1,9,9).
\]
No singleton or pair among these elements generates the order-$54$ core.
Exactly twenty-five triples do.  Every one contains $8[e_9]$; after removing
that common element, the possible index pairs are
\[
\begin{split}
&(1,3),(1,4),(1,7),(1,8),(1,11),(1,12),(1,15),(1,16),\\
&(3,7),(3,15),(3,16),(4,7),(4,15),(4,16),\\
&(7,8),(7,11),(7,12),(7,16),(8,15),(8,16),\\
&(11,15),(11,16),(12,15),(12,16),(15,16).
\end{split}
\]
This is a statement about the named presentation, not a canonical generating
set for the abstract group.

\section{Reproducibility and conclusion}
The evidence binds the exact C64, corrected C69, and C70 bytes.  Producer and
checker use separate image-count routes: concrete subgroup-poset inversion
versus Birkhoff counts and epimorphism duality.  GAP independently enumerates
all 38 target subgroups in ten types.  SymPy Smith and Hermite calculations
recover the common core, all named orders, and all twenty-five triples.  Clean
replay passes and forty-two hostile semantic mutations are rejected.  C71
therefore resolves both the local pair-intersection distribution and the two
global extrema of the complement family: its common core and its total span.

\end{document}
